% ===== PRÉAMBULE CANONIQUE — à coller VERBATIM en tête de CHAQUE .tex =====
\documentclass[11pt,a4paper]{article}
\usepackage[utf8]{inputenc}\usepackage[T1]{fontenc}\usepackage{lmodern}
\usepackage[french]{babel}
\usepackage{amsmath,amssymb,mathtools}\usepackage{icomma}
\usepackage[version=4]{mhchem}          % \ce{...} : équations & formules
\usepackage{siunitx}\sisetup{locale=FR} % \qty{}{}, \si{}, \ang{}
\usepackage{chemfig}                    % \chemfig{...} : structures développées
\usepackage{xcolor}\usepackage{colortbl}
\usepackage{tikz}\usepackage{pgfplots}\pgfplotsset{compat=1.18}
\usepackage[most]{tcolorbox}\usepackage{enumitem}\usepackage{array,booktabs}
\usepackage[a4paper,margin=2.2cm]{geometry}
\usetikzlibrary{arrows.meta,calc,angles,quotes,positioning,patterns,backgrounds,shapes.geometric}
\definecolor{primary}{RGB}{222,49,99}\definecolor{primarydark}{RGB}{150,22,55}
\definecolor{defcol}{RGB}{0,114,178}\definecolor{methcol}{RGB}{52,92,148}
\definecolor{excol}{RGB}{0,135,90}\definecolor{attncol}{RGB}{200,40,55}
\tcbset{boxbase/.style={enhanced,breakable,boxrule=0.9pt,arc=2.5pt,
  left=3mm,right=3mm,top=2.3mm,bottom=2.3mm,fonttitle=\bfseries,coltitle=white}}
% --- boîtes TUTORIEL (noms EXACTS, ne pas renommer) ---
\newtcolorbox{cle}[1][]{boxbase,colback=defcol!6,colframe=defcol,
  title={\textsc{À retenir}\ifx\relax#1\relax\relax\else\textmd{ : \itshape #1}\fi}}
\newtcolorbox{methode}[1][]{boxbase,colback=methcol!7,colframe=methcol,sharp corners,
  title={\textsc{Méthode}\ifx\relax#1\relax\relax\else\textmd{ --- \itshape #1}\fi}}
\newtcolorbox{exemplebox}[1][]{boxbase,colback=excol!5,colframe=excol,
  title={\textsc{Exemple}\ifx\relax#1\relax\relax\else\textmd{ : \itshape #1}\fi}}
\newtcolorbox{piege}[1][]{boxbase,colback=attncol!6,colframe=attncol,
  title={\textsc{Piège}\ifx\relax#1\relax\relax\else\textmd{ : \itshape #1}\fi}}
\newtcolorbox{remarque}[1][]{boxbase,colback=black!4,colframe=black!45,
  title={\textsc{Remarque}\ifx\relax#1\relax\relax\else\textmd{ : \itshape #1}\fi}}
% --- boîtes EXERCICE / CORRIGÉ (noms EXACTS, auto-comptées) ---
\newtcolorbox[auto counter]{exo}[1][]{boxbase,colback=primary!4,colframe=primary,
  title={\textsc{Exercice}~\thetcbcounter\ifx\relax#1\relax\relax\else\textmd{ --- \itshape #1}\fi}}
\newtcolorbox[auto counter]{corrige}[1][]{boxbase,colback=excol!4,colframe=excol,
  title={\textsc{Corrigé}~\thetcbcounter\ifx\relax#1\relax\relax\else\textmd{ --- \itshape #1}\fi}}
\newcommand{\R}{\mathbb{R}}\newcommand{\N}{\mathbb{N}}\newcommand{\Z}{\mathbb{Z}}\newcommand{\ds}{\displaystyle}
% (un \usepackage{fancyhdr}+en-tête est possible ; il est ignoré côté web)
% =========================================================================
\begin{document}

\begin{corrige}[une même pression dans quatre unités]
On multiplie $\qty{2}{atm}$ par chaque facteur de conversion.
\begin{enumerate}[label=\alph*)]
  \item $p=2\times101325=\qty{202650}{\pascal}$.
  \item $p=\dfrac{202650}{1000}=\qty{202.65}{\kilo\pascal}$.
  \item $p=2\times1{,}013=\qty{2.026}{bar}$.
  \item $p=2\times760=\qty{1520}{mmHg}$.
\end{enumerate}
\end{corrige}

\begin{corrige}[température et rapport de températures]
Un rapport de températures n'a de sens qu'en \textbf{kelvins}.
\begin{enumerate}[label=\alph*)]
  \item $T_1=27+273=\qty{300}{\kelvin}$.
  \item $T_2=87+273=\qty{360}{\kelvin}$.
  \item $\dfrac{T_2}{T_1}=\dfrac{360}{300}=1{,}20$ : la température absolue augmente de $20\,\%$.
  \item $\dfrac{\theta_2}{\theta_1}=\dfrac{87}{27}\approx3{,}2$. Ce nombre n'a \emph{aucun sens physique} :
        on ne peut pas former un rapport de températures en \si{\degreeCelsius}. Seul le rapport en
        kelvins ($1{,}20$) est utilisable dans les lois des gaz.
\end{enumerate}
\end{corrige}

\begin{corrige}[du volume d'une pièce au volume de gaz]
On rappelle $\qty{1}{\meter\cubed}=\qty{1000}{\litre}$.
\begin{enumerate}[label=\alph*)]
  \item $V=4\times5\times2{,}5=\qty{50}{\meter\cubed}$.
  \item $V=50\times1000=\qty{50000}{\litre}$.
  \item $V=22{,}4\times1000=\qty{22400}{\litre}$.
  \item $V=0{,}25\times1000=\qty{250}{\litre}$.
\end{enumerate}
\end{corrige}

\begin{corrige}[lecture d'un manomètre]
On divise par les facteurs de conversion (une plus grande unité $\Rightarrow$ un plus petit nombre).
\begin{enumerate}[label=\alph*)]
  \item $p=\dfrac{1520}{760}=\qty{2}{atm}$.
  \item $p=2\times1{,}013=\qty{2.026}{bar}$.
  \item $p=2\times101{,}325=\qty{202.65}{\kilo\pascal}$.
  \item $\dfrac{p}{p_{\text{atm}}}=\dfrac{2}{1}=2$ : la pression vaut \textbf{deux fois} la pression
        atmosphérique.
\end{enumerate}
\end{corrige}

\end{document}
